\documentclass{amsart}
\usepackage{amsmath,amssymb,amsthm,hyperref}
\usepackage{algorithm}
\usepackage{algpseudocode}

\newtheorem{theorem}{Theorem}
\newtheorem{lemma}{Lemma}
\newtheorem{proposition}{Proposition}
\newtheorem{corollary}{Corollary}
\theoremstyle{definition}
\newtheorem{definition}{Definition}
\newtheorem{remark}{Remark}
\newtheorem{assumption}{Assumption}

\newcommand{\Chat}{\widehat{C}}
\newcommand{\Lhat}{\widehat{L}}
\newcommand{\Nhat}{\widehat{N}}
\newcommand{\E}{\mathbb{E}}
\renewcommand{\Pr}{\mathbb{P}}

\begin{document}

\title{Measurement, estimation, and provable mapping quality in heterogeneous computing}
\maketitle

\section{Overview}\label{sec:overview}

The problem asks us to (a) define computable runtime descriptors $L_i(\tau)$ and
$N_{ij}(\tau)$, (b) produce completion-time estimators $\Chat$ with controlled
accuracy, and (c) prove that minimizing the \emph{estimated} objective yields a
mapping provably close to the optimum of the \emph{exact} objective.

The mathematical heart of the matter is the following implication, which we
prove in full generality:

\begin{quote}
\emph{If the estimator is accurate in a multiplicative (relative-error) sense
uniformly over all candidate mapping/schedule pairs, then any
mapping/schedule that exactly minimizes the estimated objective is, under the
exact objective, within a factor $\tfrac{1+\varepsilon}{1-\varepsilon}$ of
optimal, for a broad and explicitly characterized class of objectives (makespan
under precedence and weighted average completion time among them).}
\end{quote}

We then close the loop with concrete measurement and estimation procedures.
The crucial subtlety, which an earlier version of this manuscript elided, is
that the descriptors and hence the completion times \emph{depend on the chosen
mapping and schedule} through contention, queueing, and precedence. We
therefore (i) make the completion time an explicit function of the schedule via
a \emph{structural demand/load model}, (ii) estimate the schedule-independent
\emph{demand} parameters of that model from runtime probes, and (iii) propagate
the estimation error through the model to obtain uniform multiplicative accuracy
\emph{simultaneously over all schedules}, with an explicit sample-complexity
bound that carries an honest accounting of estimator bias, sample dependence,
and temporal drift.

\medskip
\noindent\textbf{Roadmap.} Section~\ref{sec:model} fixes the schedule model and
the structural completion-time model. Section~\ref{sec:objective} defines
scale-coherent objectives and verifies the canonical ones, now for the
schedule-induced finish times. Section~\ref{sec:measure} defines and computes
the descriptors. Section~\ref{sec:estimate} estimates the structural demand
parameters and propagates the error to a \emph{schedule-uniform} multiplicative
accuracy guarantee, accounting for bias, dependence, and drift.
Section~\ref{sec:transfer} proves the transfer theorem and its robust and
end-to-end corollaries.

\section{Model: schedules and structural completion times}\label{sec:model}

\subsection{Schedulable units, mappings, and schedules}

\begin{definition}[Units and precedence]\label{def:units}
Let $U=T\cup S$ be the finite set of \emph{schedulable units}: tasks
$T=\{t_1,\dots,t_k\}$ and the (finite) set $S$ of all subtasks obtained by
decomposition. The data/precedence relations form a directed acyclic graph
(DAG) $G=(U,E)$: an edge $(u,v)\in E$ means $v$ may not start before $u$
finishes and (if $u,v$ are placed on different machines) before the data
produced by $u$ has been transferred to $v$'s machine. Each unit $u$ carries an
input/output \emph{data size} $d_{uv}\ge0$ for each edge $(u,v)\in E$.
\end{definition}

\begin{definition}[Mapping and schedule]\label{def:sched}
Let $M=\{m_1,\dots,m_n\}$. A \emph{mapping} is a function $x\colon U\to M$;
write $x_u=i$ if $u$ is assigned to $m_i$. A \emph{schedule} $\sigma$ for a
mapping $x$ assigns to each unit $u$ a start time $\sigma(u)\ge0$ on machine
$x_u$ and to each cross-machine edge a transfer interval, such that
\begin{enumerate}
\item[(S1)] the units mapped to the same machine are processed under a fixed
\emph{ordering} (a total order $\prec_i$ on $\{u:x_u=i\}$); the machine's
service discipline is one of the two standard regimes fixed in advance for the
whole instance: \emph{exclusive} (a unit runs alone, so $\nu_i\equiv1$ and the
processor throttle is $\phi_i\equiv1$), or \emph{processor-sharing of a fixed
degree} (at most a hardware-fixed number $\kappa^{\mathrm{ps}}_i$ of units run
concurrently and each receives an equal share, so the throttle of a unit run
with $\nu$ co-runners is $\phi_i=1/\nu\in\{1,1/2,\dots,1/\kappa^{\mathrm{ps}}_i\}$).
In either regime the concurrency count $\nu_i$ that a unit experiences is
determined by the ordering $\prec_i$ alone (it is the position of the unit within
its sharing batch), \emph{not} by demand magnitudes;
\item[(S2)] every precedence/data constraint of $G$ is respected: $v$ starts no
earlier than $u$'s finish plus the relevant transfer time for each $(u,v)\in E$,
and each link's concurrent-transfer count $c_{ij}$ (hence $\psi_{ij}$) is the
number of transfers the routing assigns to $(i,j)$ in the same batch, again a
function of the configuration's \emph{combinatorial} routing/ordering, not of
demand magnitudes.
\end{enumerate}
A \emph{configuration} is a pair $\pi=(x,\sigma)$ satisfying (S1)--(S2) together
with the problem's hard placement constraints (co-location, exclusivity,
availability). Let $\Pi$ be the (finite, nonempty) set of admissible
configurations. We assume $\Pi\neq\emptyset$ (e.g.\ any greedy list schedule on
a feasible mapping is admissible).
\end{definition}

The finiteness of $\Pi$ as a \emph{decision} set deserves a word: start times
range over a continuum, but for any fixed mapping the relevant schedules are
determined up to combinatorial choice (the relative order of units on each
machine and the routing of transfers), and the optimal start times are then the
earliest feasible ones given that order. Thus the optimization ranges over a
finite set of \emph{orderings}, each inducing a unique left-shifted schedule;
we identify $\Pi$ with this finite set. This is the standard reduction used for
shop and DAG scheduling.

\subsection{Structural completion-time model}

The completion time of a unit is \emph{not} a primitive constant: it depends on
the machine, on the unit's intrinsic resource \emph{demand}, and on the
\emph{load} the configuration places on that machine and its links. We make
this explicit.

\begin{definition}[Demand parameters]\label{def:demand}
Each unit $u$ has a \emph{demand vector} $d_u\in\mathbb{R}^{q}_{\ge0}$
(e.g.\ floating-point operations, memory footprint, I/O volume) intrinsic to the
unit and independent of placement. Each cross-machine edge $(u,v)$ carries a
transferred data size $d_{uv}\ge0$. Collect all of these into the
\emph{demand profile} $\mathbf d=(\,(d_u)_{u},(d_{uv})_{(u,v)\in E}\,)$, a fixed
but a priori unknown vector in $\mathbb{R}^Q_{\ge0}$.
\end{definition}

\begin{definition}[Endogenous load model]\label{def:load}
The descriptors that enter the completion functions split into two parts:
\begin{enumerate}
\item[(i)] \emph{Schedule-independent structural constants} $\theta_i$ for each
machine and $(\beta_{ij},\lambda_{ij})$ for each link -- the \emph{nominal}
per-resource processing rates, nominal link bandwidth, and link latency. These
are intrinsic hardware/network parameters; they do not depend on the
configuration.
\item[(ii)] \emph{Configuration-induced contention state}: the part of $L_i,N$
that reflects \emph{who is running where and when} -- queue occupancy,
CPU/memory utilization, and the number of transfers concurrently sharing a
link. We model this explicitly as a \emph{known, computable} function of the
configuration:
\[
L_i(\pi)=\Lambda_i\big(\pi;\theta\big),\qquad
N_{ij}(\pi)=\mathcal N_{ij}\big(\pi;\beta,\lambda\big),
\]
where $\Lambda_i,\mathcal N_{ij}$ are the deterministic load maps induced by
the schedule. For instance, the realized rate on $m_i$ is the nominal rate
divided by the number $\nu_i(\pi)$ of units the schedule runs concurrently on
$m_i$ in the relevant interval (processor sharing), and a link's effective
bandwidth is its nominal bandwidth divided by the number $c_{ij}(\pi)$ of
transfers it carries concurrently (bandwidth sharing).
\end{enumerate}
The point of (ii) is that the contention state is \emph{determined by the
configuration we are evaluating}, hence \emph{computable} once the
schedule-independent constants $(\theta,\beta,\lambda)$ are known -- it need
not, and cannot, be probed per configuration. (The combinatorial quantities
$\nu_i(\pi),c_{ij}(\pi)$ depend on the \emph{schedule ordering}, not on the
demand magnitudes.)
\end{definition}

\begin{assumption}[Homogeneous contention regime]\label{ass:regime}
The realized per-unit completion time has the form
\begin{equation}\label{eq:Cform}
C_u(\pi)=g_{x_u}\big(d_u;L_{x_u}(\pi),N(\pi)\big)
=\frac{\langle\kappa_{x_u}(\theta_{x_u}),\,d_u\rangle}{\phi_{x_u}(\pi)},
\end{equation}
where $\kappa_i(\theta_i)>0$ is the nominal (uncontended) rate vector
determined by the \emph{schedule-independent} constant $\theta_i$, and
$\phi_i(\pi)\in(0,1]$ is a \emph{contention-throttle factor} determined by the
configuration's concurrency state (e.g.\ $\phi_i(\pi)=1/\nu_i(\pi)$ for
processor sharing) and \emph{independent of the demand magnitudes} $\mathbf d$.
Likewise each transfer's data-dependent part is
$d_{uv}/\big(\beta_{x_vx_u}\,\psi_{x_vx_u}(\pi)\big)$ with a demand-independent
sharing factor $\psi_{ij}(\pi)\in(0,1]$ determined by $c_{ij}(\pi)$.
\end{assumption}

\begin{remark}[Why this regime, and its scope]\label{rem:regime}
Assumption~\ref{ass:regime} is the precise content of the standard
linear-throughput-with-fair-sharing model used throughout scheduling theory:
work is linear in demand, and contention rescales the \emph{rate} by a factor
that depends on \emph{how many} competitors are present, not on \emph{how large}
each competitor's demand is. It is exactly the property that makes the realized
completion time degree-one homogeneous in $\mathbf d$ even though the load is
configuration-dependent (Lemma~\ref{lem:homog}). It does \emph{not} assume the
load is frozen or configuration-independent: the throttle $\phi_i(\pi)$ and the
descriptor $L_i(\pi)$ genuinely vary with $\pi$, resolving the central modeling
tension. Memory-hierarchy effects that make $g_i$ nonlinear in $d_u$ (e.g.\
cache thresholds) violate exact degree-one homogeneity; for those we give in
Lemma~\ref{lem:lipschitz} a Lipschitz residual bound that degrades the constant
gracefully and recovers the homogeneous result as the nonlinearity vanishes.
This is an explicit, scoped modeling commitment, not a hidden assumption.
\end{remark}

\begin{definition}[Structural completion-time / finish-time model]\label{def:struct}
Under Assumption~\ref{ass:regime}, given a configuration $\pi=(x,\sigma)$ the
\emph{realized completion time} of unit $u$ is $C_u(\pi)$ of \eqref{eq:Cform},
with the contention throttle $\phi_{x_u}(\pi)$ computed from $\pi$, and the
\emph{realized finish time} $f_u(\pi)$ is the schedule-induced completion time
of the earliest feasible left-shifted schedule, i.e.\ the unique fixed point of
the recursion
\begin{equation}\label{eq:recurdef}
f_u(\pi)=\max\Big(f_{p(u)}(\pi),\ \max_{(v,u)\in E}\big(f_v(\pi)+\theta_{vu}(\pi)\big)\Big)+C_u(\pi),
\end{equation}
where $p(u)$ is $u$'s immediate machine-predecessor under $\pi$'s ordering
($f_{p(u)}:=0$ if none) and $\theta_{vu}(\pi)=\lambda_{x_vx_u}+
d_{vu}/(\beta_{x_vx_u}\psi_{x_vx_u}(\pi))$ is the cross-edge transfer delay.
\end{definition}

\begin{lemma}[Well-posedness: the throttles are ordering-determined and the
finish-time recursion has a unique, non-circular solution]\label{lem:wellposed}
Fix a configuration $\pi=(x,\sigma)$ with its batching/ordering data of (S1)--(S2).
Then:
\begin{enumerate}
\item[\textup{(i)}] The processor throttles $\phi_i(\pi)\in(0,1]$ and the link
sharing factors $\psi_{ij}(\pi)\in(0,1]$ are determined by the
\emph{combinatorial} batching/ordering of $\pi$ alone, independently of the
demand profile $\mathbf d$ and of the realized start/finish times.
\item[\textup{(ii)}] Given these throttles, the realized completion times
$C_u(\pi)$ and transfer delays $\theta_{vu}(\pi)$ of \eqref{eq:Cform} and
Definition~\ref{def:struct} are fixed nonnegative numbers, and the finish-time
recursion \eqref{eq:recurdef} has a unique solution $(f_u(\pi))_{u\in U}$,
computable by a single topological-order pass.
\end{enumerate}
\end{lemma}

\begin{proof}
\textup{(i)} By (S1) the per-unit concurrency count $\nu_i$ is the size of the
unit's sharing batch (a position within the fixed ordering $\prec_i$ and the
fixed service degree $\kappa^{\mathrm{ps}}_i$), and by (S2) each link's
concurrent-transfer count $c_{ij}$ is the size of the corresponding routing
batch. Both are integer functions of the configuration's batching/ordering data,
which (S1)--(S2) declare in advance and which contain \emph{no} reference to the
demand magnitudes or to the continuous start times: a sharing batch is the set of
units the discipline groups together by ordinal position, not by how long they
take. Hence $\phi_i(\pi)=1/\nu_i$ and $\psi_{ij}(\pi)=1/c_{ij}$ are
$\mathbf d$-independent constants determined by $\pi$.

\textup{(ii)} With $\phi_{x_u}(\pi),\psi_{x_vx_u}(\pi)$ thus fixed, each
$C_u(\pi)=\langle\kappa_{x_u},d_u\rangle/\phi_{x_u}(\pi)\ge0$ and each
$\theta_{vu}(\pi)=\lambda_{x_vx_u}+d_{vu}/(\beta_{x_vx_u}\psi_{x_vx_u}(\pi))\ge0$
is a fixed nonnegative real that does \emph{not} depend on any $f$. Consider the
directed graph $H_\pi$ on $U$ whose arcs are (a) the DAG edges $(v,u)\in E$ and
(b) the immediate-machine-predecessor arcs $p(u)\to u$ induced by each ordering
$\prec_i$. Arcs of type (a) respect the partial order of the DAG $G$; arcs of
type (b) point from an earlier to a later unit in the total order $\prec_i$ on a
single machine. A total order extending all the $\prec_i$ and a topological order
of $G$ exists simultaneously, because $G$ is acyclic and (S2) forbids a
machine-ordering arc $p(u)\to u$ from contradicting a precedence arc $u\to p(u)$
(else $\pi$ would be infeasible, as $u$ could not both precede and start after
$p(u)$). Concretely, the start times $\sigma$ of an admissible $\pi$ furnish a
real-valued potential that is strictly increasing along every arc of $H_\pi$
(type (a): a successor starts after a predecessor finishes; type (b): a later
unit starts no earlier than an earlier one), so $H_\pi$ is acyclic. The recursion
\eqref{eq:recurdef} expresses $f_u$ as a $\max$ of sums over the in-neighbours of
$u$ in $H_\pi$ plus the constant $C_u(\pi)$; on a finite DAG such a recursion has
a unique solution obtained by evaluating the units in any topological order of
$H_\pi$ (each $f_u$ depends only on already-evaluated predecessors). This is the
earliest-feasible left-shifted schedule.
\end{proof}

\begin{remark}[No circularity between throttle and finish time]\label{rem:nocirc}
Lemma~\ref{lem:wellposed} dispels the apparent circularity (raised in review)
that ``the throttle depends on which units overlap, which depends on the finish
times, which depend on the throttle.'' Under the ordering/batching semantics of
(S1)--(S2) the overlap structure is \emph{decided combinatorially by the
configuration} (who is batched with whom), and the throttle reads off that fixed
combinatorial structure. The finish times are then computed \emph{downstream} by
a well-founded topological recursion. The optimizer ranges over the finite set of
admissible batchings/orderings $\Pi$; for each, the throttles -- and hence the
$\mathbf d$-homogeneity of completion and finish times -- are pinned down before
any finish time is computed. This is exactly the standard list-/batch-schedule
reduction, made explicit here so that Lemma~\ref{lem:homog} below rests on a
well-posed object.
\end{remark}

The key structural fact we will use is that $C_u(\pi)$ -- and through it the
finish times and every objective below -- is \emph{positively homogeneous of
degree one in the demand profile $\mathbf d$}, with the configuration held
fixed. Crucially, this holds \emph{even though the load $L_i(\pi)$ depends on
$\pi$}, because under Assumption~\ref{ass:regime} the configuration fixes the
demand-independent throttle factors $\phi_i(\pi),\psi_{ij}(\pi)$.

\begin{lemma}[Homogeneity of the finish times in demand]\label{lem:homog}
Fix a configuration $\pi\in\Pi$. For every $c>0$, replacing $\mathbf d$ by
$c\,\mathbf d$ replaces every realized completion time $C_u(\pi)$ by
$c\,C_u(\pi)$, every transfer delay's data-dependent part by $c$ times itself,
and hence every left-shifted finish time $f_u(\pi)$ by $c\,f_u(\pi)$, provided
the latency part of transfers and the schedule \emph{ordering} are held fixed.
\end{lemma}

\begin{proof}
By \eqref{eq:Cform}, $C_u(\pi)=\langle\kappa_{x_u}(\theta_{x_u}),d_u\rangle/\phi_{x_u}(\pi)$.
The throttle $\phi_{x_u}(\pi)$ and rate $\kappa_{x_u}(\theta_{x_u})$ are fixed
by the configuration (Assumption~\ref{ass:regime}: $\phi$ depends on the
concurrency count, not on demand magnitudes), so $C_u(\pi)$ is linear in $d_u$
and scaling $d_u\mapsto cd_u$ scales $C_u(\pi)\mapsto cC_u(\pi)$. A cross-edge
transfer delay is $\theta_{vu}(\pi)=\lambda_{x_vx_u}+
d_{vu}/(\beta_{x_vx_u}\psi_{x_vx_u}(\pi))$; with $\beta,\psi$ fixed by $\pi$, its
data-dependent part $d_{vu}/(\beta_{x_vx_u}\psi_{x_vx_u}(\pi))$ scales by $c$.
By Lemma~\ref{lem:wellposed} the throttles $\phi_i(\pi),\psi_{ij}(\pi)$ are
ordering-determined constants, so scaling $\mathbf d\mapsto c\mathbf d$ leaves
them unchanged; the recursion \eqref{eq:recurdef} is then a well-founded
topological evaluation (Lemma~\ref{lem:wellposed}(ii)) of a finite max of sums of
completion and transfer-delay terms along the arcs of $H_\pi$. If \emph{all}
additive constituents that depend on $\mathbf d$ scale by $c$ and the latency
constants $\lambda_{ij}$ are treated as fixed schedule overhead, then by
induction over the topological order of $H_\pi$ each $f_u$ is a max/sum of
$c$-scaled nonnegative terms, hence scales by $c$ (the base case is a source of
$H_\pi$, with $f_u=C_u$). (When
$\lambda_{ij}>0$, they form an additive $\mathbf d$-independent constant; the
homogeneity holds for the work-dependent component and
Remark~\ref{rem:latency} records the additive correction.)
\end{proof}

\section{Objectives and scale-coherence}\label{sec:objective}

We now phrase the objective directly on configurations through the realized
finish times, and isolate the structural property that drives the transfer
theorem.

\begin{definition}[Objective on configurations]\label{def:obj}
An \emph{objective} is a function $F\colon\mathbb{R}^Q_{\ge0}\times\Pi\to
\mathbb{R}_{\ge0}$, $F(\mathbf d,\pi)$, the cost of running configuration $\pi$
under demand profile $\mathbf d$. The canonical objectives are
\begin{align}
\text{(makespan)}\quad &F_{\max}(\mathbf d,\pi):=\max_{u\in U} f_u(\pi),
\label{eq:makespan}\\
\text{(weighted completion)}\quad &F_{w}(\mathbf d,\pi):=\sum_{u\in U} w_u\, f_u(\pi),
\label{eq:wsum}
\end{align}
with $w_u\ge0$, where $f_u(\pi)$ are the realized finish times of
Definition~\ref{def:struct}. Makespan is the time the last unit finishes —
correctly accounting, through $f_u$, for overlap, waiting, and precedence
(unlike the naive sum of per-unit times).
\end{definition}

\begin{definition}[Demand-scale-coherence]\label{def:coherent}
An objective $F$ is \emph{demand-scale-coherent} if it is built from the finish
times $f_u(\pi)$ by a map that is (i) \emph{monotone} (increasing each $f_u$
weakly increases $F$) and (ii) \emph{positively homogeneous of degree one in the
finish-time vector}: $F$ evaluated at finish times $(c f_u)_u$ equals
$c\cdot F$ evaluated at $(f_u)_u$, for all $c>0$. Combined with
Lemma~\ref{lem:homog}, this yields the operative property: for fixed $\pi$ and
all $c>0$,
\begin{equation}\label{eq:scale}
F(c\,\mathbf d,\pi)=c\,F(\mathbf d,\pi),
\end{equation}
and monotonicity in $\mathbf d$: $\mathbf d\le\mathbf d'$ (entrywise) implies
$F(\mathbf d,\pi)\le F(\mathbf d',\pi)$ for every $\pi$.
\end{definition}

\begin{lemma}[The canonical objectives are demand-scale-coherent]\label{lem:canonical}
$F_{\max}$ and $F_w$ are demand-scale-coherent.
\end{lemma}

\begin{proof}
By Lemma~\ref{lem:homog}, scaling $\mathbf d\mapsto c\mathbf d$ with $\pi$ fixed
scales every $f_u(\pi)\mapsto c\,f_u(\pi)$. Then
$F_{\max}=\max_u f_u$ scales as $\max_u c f_u=c\max_u f_u$, and
$F_w=\sum_u w_u f_u$ scales as $\sum_u w_u c f_u=c\sum_u w_u f_u$; both are
positively homogeneous of degree one in $(f_u)_u$, giving \eqref{eq:scale}.
Monotonicity: increasing any demand entry weakly increases every realized
completion time (each $\kappa_i>0$) and hence, through the left-shift
recursion, weakly increases every finish time; $\max$ and nonnegatively
weighted sums are monotone in their arguments. Hence both objectives are
monotone in $\mathbf d$.
\end{proof}

\begin{remark}\label{rem:classes}
Demand-scale-coherence holds for any monotone, degree-one-homogeneous aggregate
of finish times: $\ell_p$ norms of the finish-time vector ($p\ge1$), maximum
weighted finish time $\max_u w_u f_u$, makespan over per-machine loads, and
convex combinations thereof. The transfer theorem below uses \emph{only}
Definition~\ref{def:coherent} and so applies to all of these verbatim.
\end{remark}

\section{Measurement: computing descriptors and demands (part (a))}\label{sec:measure}

\begin{definition}[Descriptors]\label{def:descriptors}
At sampling instant $\tau$ each machine $m_i$ exports
$L_i(\tau)=(q_i(\tau),\rho^{\mathrm{cpu}}_i(\tau),\rho^{\mathrm{mem}}_i(\tau),
a_i(\tau),\theta_i)\in\mathbb{R}^5$ (queue length, CPU/memory utilization,
availability, nominal throughput), and each link exports
$N_{ij}(\tau)=(\beta_{ij}(\tau),\lambda_{ij}(\tau))\in\mathbb{R}^2_{\ge0}$
(effective bandwidth, latency). Consistent with Definition~\ref{def:load}, the
\emph{schedule-independent} components -- nominal throughput $\theta_i$
(equivalently the uncontended rate $\kappa_i(\theta_i)$), nominal bandwidth
$\beta_{ij}$, and latency $\lambda_{ij}$ -- are the structural constants entering
\eqref{eq:Cform}; the \emph{contention} components ($q_i,\rho^{\mathrm{cpu}}_i,
\rho^{\mathrm{mem}}_i$) are diagnostic and are reproduced \emph{predictively} for
any candidate configuration via the throttle factors $\phi_i(\pi),\psi_{ij}(\pi)$
computed from the schedule's concurrency counts, rather than measured per
configuration.
\end{definition}

\begin{algorithm}[h]
\caption{ComputeDescriptors}\label{alg:descr}
\begin{algorithmic}[1]
\Require window length $r\in\mathbb N$; raw streams $L_i(\tau_1),\dots,L_i(\tau_r)$, $N_{ij}(\tau_1),\dots,N_{ij}(\tau_r)$.
\Ensure averaged descriptors $\bar L_i,\bar N_{ij}$.
\For{$i=1$ \textbf{to} $n$}
  \State $\bar L_i\gets\frac1r\sum_{p=1}^r L_i(\tau_p)$
  \For{$j\ne i$} \State $\bar N_{ij}\gets\frac1r\sum_{p=1}^r N_{ij}(\tau_p)$ \EndFor
\EndFor
\State \Return $(\bar L_i)_i,(\bar N_{ij})_{ij}$
\end{algorithmic}
\end{algorithm}

\noindent\textbf{Correctness, termination, complexity.} A fixed number of
additions/divisions: it computes exactly the empirical means, terminates, and
runs in $O(rn^2)$ time, $O(n^2)$ space.

The descriptors feed the estimator in two ways: their schedule-independent
components instantiate the nominal rate vectors $\kappa_i(\theta_i)$ and transfer
constants $(\beta_{ij},\lambda_{ij})$ of the structural model (estimated to
relative accuracy \eqref{eq:rateacc}; Remark~\ref{rem:rates}), and the
work-counter streams are used to estimate the demand profile $\mathbf d$, as we
describe next. The contention-induced descriptor values are not probed per
configuration; they are recovered from the schedule via the throttle factors.

\section{Estimation: demands $\Rightarrow$ schedule-uniform accuracy (part (b))}\label{sec:estimate}

The earlier version probed the completion time $C_{ui}$ \emph{directly}, which is
unsound: $C_{ui}$ depends on the configuration, so a single probe corresponds to
a single load condition and cannot certify accuracy under the (different) load
of an optimal schedule, and one cannot probe a configuration without deploying
it. We instead estimate the \emph{schedule-independent} structural parameters --
the demand profile $\mathbf d$ together with the nominal rates and bandwidths
(Remark~\ref{rem:rates}) -- and \emph{compute} $\Chat$ for every configuration
via the known structural model, supplying each configuration's own,
exactly-known contention throttles. This converts a per-configuration accuracy
problem (with $|\Pi|$ exponentially large) into a per-parameter accuracy problem
(with $Q'$ entries), and propagates a uniform relative error to \emph{all}
configurations at once.

\begin{assumption}[Probing model]\label{ass:probe}
For each demand coordinate $\ell\in\{1,\dots,Q\}$ (i.e.\ each $d_u$ component or
each $d_{uv}$), we obtain measured probes $Y^{(1)}_\ell,\dots,Y^{(r_\ell)}_\ell$
of $d_\ell$ — e.g.\ instruction/FLOP counters, allocated-byte counters, and
transferred-byte counters, read under a lightly loaded calibration run or
directly from hardware performance counters that measure \emph{work}, not
elapsed time. We assume:
\begin{enumerate}
\item[(M1$'$)] \emph{Bounded bias.} $|\E[\bar Y_\ell]-d_\ell|\le\beta\,d_\ell$
for a known $\beta\in[0,1)$, where $\bar Y_\ell=\frac1{r_\ell}\sum_p Y^{(p)}_\ell$.
(Counters of \emph{work performed} are nearly unbiased; $\beta$ budgets any
residual systematic error. Setting $\beta=0$ recovers exact unbiasedness.)
\item[(M2)] \emph{Boundedness.} $0\le Y^{(p)}_\ell\le b_\ell$ a.s., $b_\ell>0$ known.
\item[(M3)] \emph{Positivity.} $d_\ell\ge a_\ell>0$ known.
\item[(M4)] \emph{Finite-range dependence (block structure).} The probe stream is
\emph{$m$-dependent}: probes more than $m$ steps apart are independent (e.g.\
because each microbenchmark fully completes and resets shared state before the
next begins, with $m$ accounting for residual cache/pipeline carry-over). We
take probes in disjoint blocks: partition the $r_\ell$ probes into
$B_\ell:=\lfloor r_\ell/(m+1)\rfloor$ groups separated by gaps of $m$ discarded
probes, so the per-block averages
$Z^{(1)}_\ell,\dots,Z^{(B_\ell)}_\ell$ (each the mean of one block's retained
probes) are \emph{independent}, identically distributed, bounded in $[0,b_\ell]$,
with common mean $\mu_\ell=\E[\bar Y_\ell]$. The \emph{effective sample size} is
the number of independent blocks, $r_\ell^{\mathrm{eff}}:=B_\ell\ge
r_\ell/(2(m+1))$. (For i.i.d.\ micro-probes, $m=0$ and
$r_\ell^{\mathrm{eff}}=r_\ell$.) We use $\bar Y_\ell:=\frac1{B_\ell}\sum_{b}Z^{(b)}_\ell$
as the estimator.
\end{enumerate}
\end{assumption}

\begin{proposition}[Per-coordinate multiplicative accuracy]\label{prop:samples}
Fix $\varepsilon\in(0,1)$ with $\varepsilon>\beta$, and $\delta\in(0,1)$. Put
$\varepsilon_0:=\varepsilon-\beta>0$. Under
Assumption~\ref{ass:probe}, if
\begin{equation}\label{eq:rbound}
r_\ell^{\mathrm{eff}}\ \ge\ \frac{b_\ell^2}{2\,a_\ell^2\,\varepsilon_0^2}\,
\ln\frac2\delta
\qquad\Big(\text{i.e. } r_\ell\ge 2(m+1)\,\frac{b_\ell^2}{2a_\ell^2\varepsilon_0^2}\ln\tfrac2\delta\Big),
\end{equation}
then $\Pr\big[(1-\varepsilon)d_\ell\le \bar Y_\ell\le(1+\varepsilon)d_\ell\big]\ge1-\delta$.
\end{proposition}

\begin{proof}
Write $\mu_\ell:=\E[\bar Y_\ell]$. By (M4) the estimator $\bar Y_\ell$ is the
mean of $B_\ell=r_\ell^{\mathrm{eff}}$ \emph{independent} block averages
$Z^{(b)}_\ell\in[0,b_\ell]$ with common mean $\mu_\ell$. Hoeffding's inequality
for bounded independent variables applies verbatim (no heuristic substitution):
for any $s>0$,
\[
\Pr\big[|\bar Y_\ell-\mu_\ell|\ge s\big]\le 2\exp\!\Big(-\frac{2 r_\ell^{\mathrm{eff}} s^2}{b_\ell^2}\Big).
\]
Take $s=\varepsilon_0 a_\ell$. The right side is $\le\delta$ exactly when
\eqref{eq:rbound} holds. On the complementary event,
$|\bar Y_\ell-\mu_\ell|<\varepsilon_0 a_\ell\le\varepsilon_0 d_\ell$ by (M3).
Combining with the bias bound (M1$'$), $|\mu_\ell-d_\ell|\le\beta d_\ell$, the
triangle inequality gives
\[
|\bar Y_\ell-d_\ell|\ \le\ |\bar Y_\ell-\mu_\ell|+|\mu_\ell-d_\ell|
\ <\ \varepsilon_0 d_\ell+\beta d_\ell\ =\ \varepsilon\, d_\ell,
\]
i.e.\ $(1-\varepsilon)d_\ell<\bar Y_\ell<(1+\varepsilon)d_\ell$.
\end{proof}

\begin{corollary}[Uniform demand accuracy]\label{cor:uniform}
Define $\widehat{\mathbf d}:=(\bar Y_\ell)_{\ell=1}^Q$. If for each coordinate
$r_\ell^{\mathrm{eff}}\ge\frac{b_\ell^2}{2a_\ell^2\varepsilon_0^2}\ln\frac{2Q}{\delta}$,
then with probability $\ge1-\delta$,
\begin{equation}\label{eq:Deps}
(1-\varepsilon)\,d_\ell\ \le\ \widehat d_\ell\ \le\ (1+\varepsilon)\,d_\ell
\qquad\text{for all }\ell=1,\dots,Q.
\end{equation}
\end{corollary}

\begin{proof}
Apply Proposition~\ref{prop:samples} with $\delta\mapsto\delta/Q$ to each
coordinate and take a union bound over the $Q$ coordinates.
\end{proof}

We call event \eqref{eq:Deps} the $\varepsilon$-accuracy event $E_\varepsilon$.
The expected relative error per coordinate is
$\E|\bar Y_\ell-d_\ell|/d_\ell\le\beta+\sqrt{\operatorname{Var}(\bar Y_\ell)}/d_\ell
\le\beta+b_\ell/(2a_\ell\sqrt{r_\ell^{\mathrm{eff}}})=\beta+O(1/\sqrt r)$,
so for $\beta\to0$ it is $O(1/\sqrt r)$.

\begin{lemma}[Error propagation: demands and rates $\Rightarrow$ completion times]\label{lem:propagate}
Suppose the accuracy event $E_\varepsilon$ of \eqref{eq:Deps} holds for the
demand profile $\mathbf d$ and that, in addition, the estimated nominal
structural constants are relatively accurate at the same level, i.e.\ the
nominal rate vectors and bandwidths satisfy
\begin{equation}\label{eq:rateacc}
(1-\varepsilon)\,\kappa_{i\ell}\le\widehat\kappa_{i\ell}\le(1+\varepsilon)\,\kappa_{i\ell},
\qquad
(1-\varepsilon)\,\beta_{ij}\le\widehat\beta_{ij}\le(1+\varepsilon)\,\beta_{ij}
\end{equation}
for all $i,j,\ell$ (these are estimated by the same procedure of
Section~\ref{sec:estimate} applied to the descriptor streams; see
Remark~\ref{rem:rates}). Build the estimated completion and transfer times by
evaluating the \emph{same} structural model \eqref{eq:Cform} at the estimated
constants \emph{and the exactly-computed, demand-independent throttle factors of
the configuration being evaluated}:
\[
\Chat_u(\pi):=\frac{\langle\widehat\kappa_{x_u},\widehat d_u\rangle}{\phi_{x_u}(\pi)},
\qquad
\widehat\theta_{vu}(\pi):=\lambda_{x_vx_u}+\frac{\widehat d_{vu}}
{\widehat\beta_{x_vx_u}\,\psi_{x_vx_u}(\pi)} .
\]
Here $\phi_{x_u}(\pi)$ and $\psi_{x_vx_u}(\pi)$ are the \emph{same} throttle and
sharing factors that appear in the realized $C_u(\pi),\theta_{vu}(\pi)$, because
they are combinatorial functions of the schedule (concurrency counts) and are
known exactly for the configuration under evaluation -- they are \emph{not}
estimated. Set $\varepsilon_\star:=\tfrac{2\varepsilon}{1-\varepsilon}\ge2\varepsilon$.
Then for every configuration $\pi\in\Pi$ and every unit $u$,
\begin{equation}\label{eq:Ceps}
(1-\varepsilon_\star)\,C_u(\pi)\ \le\ \Chat_u(\pi)\ \le\ (1+\varepsilon_\star)\,C_u(\pi),
\end{equation}
and the same two-sided multiplicative bound holds for every realized finish time
$\widehat f_u(\pi)$ versus $f_u(\pi)$ (work-dependent components; see
Remark~\ref{rem:latency} for latencies).
\end{lemma}

\begin{proof}
Fix $\pi$. By \eqref{eq:Cform} the truth and estimate share the \emph{same}
exactly-known throttle $\phi:=\phi_{x_u}(\pi)\in(0,1]$:
\[
C_u(\pi)=\frac{\langle\kappa_{x_u},d_u\rangle}{\phi},
\qquad
\Chat_u(\pi)=\frac{\langle\widehat\kappa_{x_u},\widehat d_u\rangle}{\phi}.
\]
The factor $1/\phi>0$ is common and cancels in every ratio, so it suffices to
compare the nominal works $W_u:=\langle\kappa_{x_u},d_u\rangle$ and
$\widehat W_u:=\langle\widehat\kappa_{x_u},\widehat d_u\rangle$. For each
coordinate $\ell$, $\kappa_{x_u\ell}d_{u\ell}\ge0$ and, by \eqref{eq:Deps} and
\eqref{eq:rateacc},
\[
(1-\varepsilon)^2\,\kappa_{x_u\ell}\,d_{u\ell}
\ \le\ \widehat\kappa_{x_u\ell}\,\widehat d_{u\ell}
\ \le\ (1+\varepsilon)^2\,\kappa_{x_u\ell}\,d_{u\ell}.
\]
Summing the nonnegative coordinatewise products gives
$(1-\varepsilon)^2 W_u\le\widehat W_u\le(1+\varepsilon)^2 W_u$, hence
$(1-\varepsilon)^2 C_u(\pi)\le\Chat_u(\pi)\le(1+\varepsilon)^2 C_u(\pi)$. Now
$(1+\varepsilon)^2=1+2\varepsilon+\varepsilon^2\le1+\frac{2\varepsilon}{1-\varepsilon}
=1+\varepsilon_\star$ (since $2\varepsilon+\varepsilon^2\le\frac{2\varepsilon}{1-\varepsilon}$
for $\varepsilon\in(0,1)$) and $(1-\varepsilon)^2=1-2\varepsilon+\varepsilon^2
\ge1-2\varepsilon\ge1-\varepsilon_\star$, so \eqref{eq:Ceps} follows. The same
argument applied to each transfer's data-dependent part
$d_{vu}/(\beta_{x_vx_u}\psi_{x_vx_u}(\pi))$ -- demand in the numerator,
bandwidth in the denominator, sharing factor $\psi(\pi)$ common and exactly
known -- gives the work part of the transfer delay sandwiched in
$(1\pm\varepsilon_\star)$ of its true value: indeed
$\widehat d_{vu}/\widehat\beta\in[(1-\varepsilon)/(1+\varepsilon),
(1+\varepsilon)/(1-\varepsilon)]\cdot(d_{vu}/\beta)\subseteq
[1-\varepsilon_\star,1+\varepsilon_\star]\cdot(d_{vu}/\beta)$, using
$\frac{1+\varepsilon}{1-\varepsilon}=1+\varepsilon_\star$ and
$\frac{1-\varepsilon}{1+\varepsilon}\ge1-\varepsilon_\star$. (The latency
$\lambda_{ij}$ is an additive, $\mathbf d$-independent, exactly known constant;
it is handled in Remark~\ref{rem:latency}.)

We now prove the finish-time bound by induction on the topological order of the
DAG $G$ (the schedule ordering of $\pi$ is fixed throughout, and with it all
throttle/sharing factors). Recall the recursion \eqref{eq:recurdef}: writing
$p(u)$ for the immediate machine-predecessor of $u$ under $\pi$ (with
$f_{p(u)}:=0$ if none),
\begin{equation}\label{eq:recur}
f_u(\pi)=\max\Big(f_{p(u)}(\pi),\ \max_{(v,u)\in E}\big(f_v(\pi)+\theta_{vu}(\pi)\big)\Big)+C_u(\pi),
\end{equation}
and identically for $\widehat f_u$ with $\widehat C_u,\widehat\theta$ in place of
$C_u,\theta$. We claim $(1-\varepsilon_\star)f_u(\pi)\le\widehat f_u(\pi)\le
(1+\varepsilon_\star)f_u(\pi)$ for every $u$. \emph{Base case:} a source $u$ (no
predecessors) has $f_u=C_u$ and $\widehat f_u=\widehat C_u$, so the claim is
\eqref{eq:Ceps}. \emph{Inductive step:} assume the claim for all units preceding
$u$ in \eqref{eq:recur}. Every term inside the outer $\max$ of \eqref{eq:recur}
is nonnegative, and by the inductive hypothesis and the term bounds just
established, each estimated term is sandwiched by $(1\pm\varepsilon_\star)$ times
its exact counterpart:
$\widehat f_{p(u)}\in(1\pm\varepsilon_\star)f_{p(u)}$, and for each $(v,u)\in E$,
$\widehat f_v+\widehat\theta_{vu}\le(1+\varepsilon_\star)f_v+(1+\varepsilon_\star)\theta_{vu}
=(1+\varepsilon_\star)(f_v+\theta_{vu})$ and symmetrically
$\ge(1-\varepsilon_\star)(f_v+\theta_{vu})$.
We use the elementary \emph{scaling-closure} facts, valid for finite families of
nonnegative reals $\{a_r\}$ with $(1-\eta)a_r\le\hat a_r\le(1+\eta)a_r$:
$\max_r\hat a_r\in[(1-\eta)\max_r a_r,(1+\eta)\max_r a_r]$ and
$\sum_r\hat a_r\in[(1-\eta)\sum_r a_r,(1+\eta)\sum_r a_r]$ (each is immediate
from the monotonicity of $\max$ and $\sum$ in nonnegative arguments and the
common factor). Applying the $\max$-closure (with $\eta=\varepsilon_\star$) to
the outer maximum and then the sum-closure to add $C_u$ (itself
$(1\pm\varepsilon_\star)$-bounded by \eqref{eq:Ceps}) yields
$\widehat f_u(\pi)\in[(1-\varepsilon_\star)f_u(\pi),(1+\varepsilon_\star)f_u(\pi)]$.
This completes the induction and the lemma.
\end{proof}

\begin{remark}[Estimating the nominal rates and bandwidths]\label{rem:rates}
The relative accuracy \eqref{eq:rateacc} is obtained by the identical
estimation machinery of Section~\ref{sec:estimate} applied to descriptor
probes: the nominal rate $\kappa_{i\ell}$ is read from uncontended calibration
microbenchmarks ($\nu_i=1$, so the measured elapsed time directly reports
$1/\kappa_{i\ell}$ per unit work) and the nominal bandwidth $\beta_{ij}$ from
uncontended transfer probes ($c_{ij}=1$). These are schedule-independent
hardware constants, so a single calibration -- not one per configuration --
suffices; Proposition~\ref{prop:samples} and Corollary~\ref{cor:uniform} apply
verbatim with these quantities adjoined to the parameter vector (the union
bound's $Q$ becomes $Q'=Q+\sum_i q+|\{\text{links}\}|$, still polynomial).
Because \eqref{eq:rateacc} is at the same relative level $\varepsilon$ as
\eqref{eq:Deps}, the combined propagation constant is $\varepsilon_\star$; if one
prefers a single budget, run all estimators at level $\varepsilon/2$, so that
$(1\pm\varepsilon/2)^2$ lands inside $(1\pm\varepsilon)$ and \eqref{eq:Ceps}
holds with $\varepsilon$ in place of $\varepsilon_\star$.
\end{remark}

\begin{lemma}[Residual bound for non-homogeneous completion functions]\label{lem:lipschitz}
Suppose the true per-unit completion function deviates from the homogeneous
model \eqref{eq:Cform} by a bounded multiplicative residual: there is
$\xi\in[0,1)$ with
\[
(1-\xi)\,\frac{\langle\kappa_{x_u},d_u\rangle}{\phi_{x_u}(\pi)}
\ \le\ C_u^{\mathrm{true}}(\pi)\ \le\
(1+\xi)\,\frac{\langle\kappa_{x_u},d_u\rangle}{\phi_{x_u}(\pi)}
\qquad\text{for all }u,\pi,
\]
i.e.\ nonlinearities (cache thresholds, memory-hierarchy effects) perturb the
linear-throughput value by at most a factor $1\pm\xi$. Then, building $\Chat$
from the homogeneous model at the estimated parameters as in
Lemma~\ref{lem:propagate}, the realized finish times satisfy, for every $\pi$
and $u$,
\[
\frac{1-\varepsilon_\star}{1+\xi}\,f_u^{\mathrm{true}}(\pi)
\ \le\ \widehat f_u(\pi)\ \le\
\frac{1+\varepsilon_\star}{1-\xi}\,f_u^{\mathrm{true}}(\pi),
\]
and consequently every demand-scale-coherent objective obeys the sandwich
\eqref{eq:sandwich} with $\eta$ replaced by
$\eta_\xi:=\frac{(1+\varepsilon_\star)(1+\xi)}{(1-\varepsilon_\star)(1-\xi)}-1$
on the upper side and the symmetric lower bound, giving the transfer ratio
\[
\frac{F(\hat\pi)}{F(\pi^\star)}\ \le\
\frac{1+\varepsilon_\star}{1-\varepsilon_\star}\cdot\frac{1+\xi}{1-\xi}.
\]
As $\xi\to0$ this recovers Theorem~\ref{thm:main} exactly; the model nonlinearity
enters only through the explicit, multiplicative factor $\frac{1+\xi}{1-\xi}$.
\end{lemma}

\begin{proof}
Write $C_u^{\mathrm{hom}}(\pi):=\langle\kappa_{x_u},d_u\rangle/\phi_{x_u}(\pi)$
for the homogeneous value. By hypothesis
$C_u^{\mathrm{true}}\in[(1-\xi),(1+\xi)]\,C_u^{\mathrm{hom}}$, and by
Lemma~\ref{lem:propagate} $\Chat_u\in[(1-\varepsilon_\star),(1+\varepsilon_\star)]
\,C_u^{\mathrm{hom}}$. Hence
$\Chat_u/C_u^{\mathrm{true}}\in[\frac{1-\varepsilon_\star}{1+\xi},
\frac{1+\varepsilon_\star}{1-\xi}]$, and the same per-unit two-sided factor
propagates to the finish times by the identical max/sum scaling-closure
induction of Lemma~\ref{lem:propagate} (the closure facts hold for any common
two-sided factor on nonnegative terms; the true finish times $f_u^{\mathrm{true}}$
satisfy the recursion \eqref{eq:recur} with $C_u^{\mathrm{true}}$ and the same
exactly-known throttle/sharing factors). Feeding these finish-time factors into
the monotone, degree-one-homogeneous objective (Definition~\ref{def:coherent})
gives the stated objective sandwich, and the transfer argument of
Theorem~\ref{thm:main} then yields the ratio.
\end{proof}

The decisive gain over direct probing: \eqref{eq:Ceps} holds for \emph{all}
$\pi\in\Pi$ \emph{simultaneously} from the \emph{single} demand-accuracy event
$E_\varepsilon$, because the structural model carries the configuration
dependence deterministically. The sample budget is $O(Q\,\varepsilon_0^{-2}
\ln(Q/\delta))$, polynomial in the input size, with no dependence on $|\Pi|$.

\begin{remark}[Temporal drift as a bias term]\label{rem:drift}
Assumption~\ref{ass:probe} (M4) requires stationarity \emph{within} the probing
horizon. A genuinely drifting parameter is incompatible with strict
stationarity, so drift is correctly modeled not by weakening (M4) but as an
\emph{additional bias} of the (still-stationarity-analyzed) block estimator
relative to the target value at decision time. Concretely, suppose the
parameter is Lipschitz-in-time with constant $\zeta_\ell$,
$|d_\ell(\tau)-d_\ell(\tau')|\le\zeta_\ell|\tau-\tau'|$, the probing horizon
spans $H$ time units, and the decision is taken at the end of the horizon, so the
target is $d_\ell:=d_\ell(\tau_{\mathrm{dec}})$. Run the concentration analysis
treating each block as drawn from the \emph{frozen} law at its own time stamp;
within a block of duration $\le H$ the law shifts by at most $\zeta_\ell H$, and
hence the estimator's mean over the horizon obeys
$|\E[\bar Y_\ell]-d_\ell|\le\beta d_\ell+\zeta_\ell H$, i.e.\ a per-coordinate
relative-bias budget
\[
\beta'_\ell:=\beta+\frac{\zeta_\ell H}{d_\ell}\ \le\ \beta+\frac{\zeta_\ell H}{a_\ell},
\]
using (M3) ($d_\ell\ge a_\ell$). Substituting $\beta'_\ell$ for $\beta$ in
Proposition~\ref{prop:samples} (so $\varepsilon_0=\varepsilon-\beta'_\ell$)
preserves every conclusion verbatim, provided one keeps
$\varepsilon>\beta'_\ell$, i.e.\ the horizon is short enough that
$\zeta_\ell H\le(\varepsilon-\beta)\,a_\ell$ \emph{per coordinate}. (The earlier
draft used a global $a_{\min}$; the correct, slightly stronger statement is the
per-coordinate one above, which we adopt.) Thus drift is handled by a
per-coordinate horizon constraint folded into the same bias accounting, with no
appeal to stationarity across the drift.
\end{remark}

\section{Transfer theorem: accuracy $\Rightarrow$ mapping quality (part (c))}\label{sec:transfer}

We argue \emph{deterministically} on $E_\varepsilon$; the high-probability
statement follows by Corollary~\ref{cor:uniform}.

For a demand-scale-coherent $F$ define the \emph{estimated objective}
$\widehat F(\pi):=F(\widehat{\mathbf d},\pi)$ — i.e.\ $F$ evaluated with the
estimated finish times $\widehat f_u(\pi)$ — and the \emph{exact objective}
$F(\pi):=F(\mathbf d,\pi)$.

\begin{lemma}[Pointwise objective sandwich]\label{lem:sandwich}
Let $F$ be demand-scale-coherent and suppose the accuracy events $E_\varepsilon$
\eqref{eq:Deps} and the rate accuracy \eqref{eq:rateacc} hold,
$\varepsilon\in(0,1)$, with propagated finish-time accuracy
$\varepsilon_\star=\frac{2\varepsilon}{1-\varepsilon}$ from
Lemma~\ref{lem:propagate} (assume $\varepsilon_\star<1$, i.e.\ $\varepsilon<1/3$,
or run estimators at level $\varepsilon/2$ as in Remark~\ref{rem:rates} so the
propagated level is $\varepsilon<1$). Then for every configuration $\pi\in\Pi$,
\begin{equation}\label{eq:sandwich}
(1-\varepsilon_\star)\,F(\pi)\ \le\ \widehat F(\pi)\ \le\ (1+\varepsilon_\star)\,F(\pi).
\end{equation}
\end{lemma}

\begin{proof}
By Lemma~\ref{lem:propagate}, $\widehat f_u(\pi)\le(1+\varepsilon_\star)f_u(\pi)$
for all $u$. By monotonicity of $F$ in the finish times and then degree-one
homogeneity (Definition~\ref{def:coherent}),
\[
\widehat F(\pi)=F\big((\widehat f_u(\pi))_u\big)\le F\big(((1+\varepsilon_\star)f_u(\pi))_u\big)
=(1+\varepsilon_\star)F(\pi).
\]
Symmetrically $\widehat f_u(\pi)\ge(1-\varepsilon_\star)f_u(\pi)$ and
$1-\varepsilon_\star>0$ give $\widehat F(\pi)\ge(1-\varepsilon_\star)F(\pi)$.
\end{proof}

\begin{remark}[A single accuracy budget]\label{rem:budget}
To avoid carrying both $\varepsilon$ and $\varepsilon_\star$, fix the desired
\emph{finish-time} accuracy $\eta\in(0,1)$ and run every estimator
(demands, rates, bandwidths) at relative level
$\varepsilon=\varepsilon(\eta):=\eta/(2+\eta)$, so that
$\varepsilon_\star=\frac{2\varepsilon}{1-\varepsilon}=\eta$. Then \eqref{eq:Ceps}
and \eqref{eq:sandwich} hold with $\eta$ in place of $\varepsilon_\star$, and the
transfer ratio below is $\frac{1+\eta}{1-\eta}$. We adopt this convention
hereafter: \emph{$\eta$ denotes the propagated, uniform-over-$\Pi$ relative
accuracy of completion and finish times}, achievable with the sample budget of
Corollary~\ref{cor:uniform} at level $\varepsilon(\eta)$.
\end{remark}

\begin{theorem}[Approximation guarantee under estimated optimization]\label{thm:main}
Let $F$ be demand-scale-coherent and suppose the finish-time accuracy
\eqref{eq:sandwich} holds at propagated level $\eta\in(0,1)$
(Remark~\ref{rem:budget}). Let $\hat\pi\in\arg\min_{\pi\in\Pi}\widehat F(\pi)$ and
$\pi^\star\in\arg\min_{\pi\in\Pi}F(\pi)$. Then
\begin{equation}\label{eq:approx}
F(\hat\pi)\ \le\ \frac{1+\eta}{1-\eta}\;F(\pi^\star).
\end{equation}
\end{theorem}

\begin{proof}
Both arg-min sets are nonempty since $\Pi$ is finite, nonempty. By the left
inequality of \eqref{eq:sandwich} at $\hat\pi$, optimality of $\hat\pi$ for
$\widehat F$ (so $\widehat F(\hat\pi)\le\widehat F(\pi^\star)$), and the right
inequality at $\pi^\star$,
\[
(1-\eta)F(\hat\pi)\le\widehat F(\hat\pi)\le\widehat F(\pi^\star)
\le(1+\eta)F(\pi^\star).
\]
Divide by $1-\eta>0$.
\end{proof}

\begin{theorem}[Robustness to approximate optimization]\label{thm:approxalg}
Let $F$ be demand-scale-coherent, the finish-time accuracy \eqref{eq:sandwich}
hold at level $\eta\in(0,1)$, and $\gamma\ge1$. If $\tilde\pi\in\Pi$ satisfies
$\widehat F(\tilde\pi)\le\gamma\min_{\pi}\widehat F(\pi)$, then
\begin{equation}\label{eq:approxalg}
F(\tilde\pi)\ \le\ \gamma\,\frac{1+\eta}{1-\eta}\;\min_{\pi\in\Pi}F(\pi).
\end{equation}
\end{theorem}

\begin{proof}
With $\pi^\star\in\arg\min_\pi F(\pi)$, chain the left inequality of
\eqref{eq:sandwich} at $\tilde\pi$, the $\gamma$-approximation hypothesis,
$\min_\pi\widehat F\le\widehat F(\pi^\star)$, and the right inequality at
$\pi^\star$:
\[
(1-\eta)F(\tilde\pi)\le\widehat F(\tilde\pi)\le\gamma\min_\pi\widehat F
\le\gamma\widehat F(\pi^\star)\le\gamma(1+\eta)F(\pi^\star).
\]
Divide by $1-\eta$. Taking $\gamma=1$ recovers Theorem~\ref{thm:main}.
\end{proof}

This is the relevant form in practice. Optimal DAG scheduling is NP-hard, so the
guarantee is stated relative to a $\gamma$-approximation \emph{oracle} for the
estimated objective; the transfer theorem is black-box in $\gamma$ and contributes
the multiplicative factor $\frac{1+\eta}{1-\eta}$ on top of whatever $\gamma$ the
oracle achieves. Concrete values of $\gamma$ are inherited from the scheduling
literature \emph{for the matching machine/communication model}: $\gamma=2-\tfrac1n$
for makespan on identical machines without communication (Graham list scheduling);
$\gamma=O(\log n)$-type or model-specific constants for unrelated machines
($R\,|\,|\,C_{\max}$, $2$-approximable via the Lenstra--Shmoys--Tardos rounding)
and for related/heterogeneous variants. We do \emph{not} claim a universal
constant $\gamma$ for the fully general DAG-with-communication problem
($P,R\,|\,\mathrm{prec},c_{ij}\,|\,F$), where no constant-factor polynomial
approximation is known in general; for that case the statement is the honest
conditional one: \emph{whatever} $\gamma$ a polynomial oracle attains on the
estimated instance, the realized cost is within $\gamma\frac{1+\eta}{1-\eta}$ of
the true optimum (Theorem~\ref{thm:approxalg}). The contribution of this paper is
precisely the $\frac{1+\eta}{1-\eta}$ \emph{estimation} factor, decoupled from
the $\gamma$ \emph{optimization} factor.

\begin{remark}[Additive latency correction]\label{rem:latency}
If nonzero pure latencies $\lambda_{ij}>0$ appear, write
$f_u(\pi)=f_u^{\mathrm{work}}(\pi)+\Lambda_u(\pi)$ where $\Lambda_u(\pi)$ is the
total accumulated latency along the critical chain to $u$ (a
$\mathbf d$-independent, exactly known constant for fixed $\pi$). Then
$\widehat f_u(\pi)$ satisfies $|\widehat f_u-f_u|\le\varepsilon f_u^{\mathrm{work}}
\le\varepsilon f_u$, so \eqref{eq:sandwich} holds verbatim; if instead one wants
the latency excluded from the relative budget, replace $\varepsilon$ by
$\varepsilon\,\sup_\pi f_u^{\mathrm{work}}(\pi)/f_u(\pi)\le\varepsilon$ — the
bound only improves. Hence latencies never worsen the ratio.
\end{remark}

\begin{corollary}[End-to-end high-probability guarantee]\label{cor:e2e}
Fix a target finish-time accuracy $\eta\in(0,1)$, a bias bound
$\beta\in[0,\varepsilon(\eta))$ with $\varepsilon(\eta)=\eta/(2+\eta)$,
$\delta\in(0,1)$, a demand-scale-coherent $F$, and a $\gamma$-approximation
oracle for $\min_\pi\widehat F$. Run \textsc{ComputeDescriptors}; estimate every
structural parameter (the $Q'$ demand, rate, and bandwidth coordinates of
Remark~\ref{rem:rates}) at relative level $\varepsilon(\eta)$ with
$r_\ell^{\mathrm{eff}}\ge\frac{b_\ell^2}{2a_\ell^2(\varepsilon(\eta)-\beta)^2}\ln\frac{2Q'}{\delta}$
probes per coordinate (Assumption~\ref{ass:probe}); and output $\tilde\pi$ from
the oracle on $\widehat F$. Then with probability $\ge1-\delta$,
\[
F(\tilde\pi)\ \le\ \gamma\,\frac{1+\eta}{1-\eta}\,\min_{\pi\in\Pi}F(\pi).
\]
\end{corollary}

\begin{proof}
By Corollary~\ref{cor:uniform} applied to all $Q'$ parameters with target
$\delta$, the accuracy events \eqref{eq:Deps} and \eqref{eq:rateacc} hold
simultaneously at level $\varepsilon(\eta)$ w.p.\ $\ge1-\delta$; on this event
Lemma~\ref{lem:propagate} gives uniform completion/finish-time accuracy at the
propagated level $\varepsilon_\star=\eta$ over \emph{all} $\pi$
(Remark~\ref{rem:budget}), and Theorem~\ref{thm:approxalg} gives the bound.
\end{proof}

\begin{remark}[Tightness]\label{rem:tight}
The factor $\frac{1+\eta}{1-\eta}$ in $\eta$ (the propagated, uniform
completion/finish-time relative accuracy of Remark~\ref{rem:budget}) is
unimprovable for demand-scale-coherent objectives, and is attained by a single
consistent estimate. Take a single unit $u$, two machines, $F=F_{\max}$ (so $f_u=C_u$, no
precedence), and \emph{two} demand coordinates $\mathbf d=(d_1,d_2)$. Let the
machines have disjoint rate vectors $\kappa_1=(1,0)$ and $\kappa_2=(0,1)$, so the
true completion times are $C_{u1}=\langle\kappa_1,\mathbf d\rangle=d_1$ and
$C_{u2}=\langle\kappa_2,\mathbf d\rangle=d_2$. Choose the exact demands
$d_2=1$ and $d_1=\frac{1+\eta}{1-\eta}\,d_2>d_2$, so machine~$2$ is
the unique true optimum with $F(\pi^\star)=d_2=1$.

Now consider the single estimate
\[
\widehat{\mathbf d}=\big((1-\eta)\,d_1,\ (1+\eta)\,d_2\big),
\]
which is one consistent vector and satisfies the accuracy event \eqref{eq:Deps}
coordinatewise: $\widehat d_1=(1-\eta)d_1$ and
$\widehat d_2=(1+\eta)d_2$ both lie in $[(1-\eta)d_\ell,
(1+\eta)d_\ell]$. The estimated completion times are
$\widehat C_{u1}=\widehat d_1=(1-\eta)d_1
=(1-\eta)\frac{1+\eta}{1-\eta}d_2=(1+\eta)d_2$ and
$\widehat C_{u2}=\widehat d_2=(1+\eta)d_2$. Thus
$\widehat C_{u1}=\widehat C_{u2}$: the two estimated objectives coincide, so an
estimated-optimal $\hat\pi$ may legitimately place $u$ on machine~$1$, giving
true cost $F(\hat\pi)=C_{u1}=d_1=\frac{1+\eta}{1-\eta}$. Hence
\[
\frac{F(\hat\pi)}{F(\pi^\star)}=\frac{1+\eta}{1-\eta},
\]
attained by an estimate obeying \eqref{eq:Deps}, so the bound of
Theorem~\ref{thm:main} is tight. (To break the tie strictly one may perturb
$d_1$ to $\frac{1+\eta}{1-\eta}d_2-\rho$ for arbitrarily small
$\rho>0$, forcing the optimizer onto machine~$1$ and approaching the ratio from
below.)
\end{remark}

\begin{algorithm}[h]
\caption{EstimatedOptimalMapping}\label{alg:map}
\begin{algorithmic}[1]
\Require target finish-time accuracy $\eta\in(0,1)$; bias bound $\beta\in[0,\varepsilon(\eta))$ with $\varepsilon(\eta)=\eta/(2+\eta)$; $\delta\in(0,1)$; demand-scale-coherent $F$; bounds $a_\ell,b_\ell$; dependence range $m$; $\gamma$-oracle; configuration family $\Pi$.
\Ensure $\tilde\pi$ with $F(\tilde\pi)\le\gamma\frac{1+\eta}{1-\eta}\min_\pi F(\pi)$ w.p.\ $\ge1-\delta$.
\State $\varepsilon\gets\varepsilon(\eta)$; $(\bar L,\bar N)\gets$ \textsc{ComputeDescriptors}
\For{each structural coordinate $\ell=1,\dots,Q'$ (demands, nominal rates $\kappa$, nominal bandwidths $\beta_{ij}$)}
  \State $r_\ell\gets\big\lceil 2(m+1)\frac{b_\ell^2}{2a_\ell^2(\varepsilon-\beta)^2}\ln\frac{2Q'}{\delta}\big\rceil$ \Comment{$2(m+1)$ converts blocks to raw probes}
  \State draw $r_\ell$ probes per Assumption~\ref{ass:probe}, form $B_\ell$ independent block means, set $\widehat\theta_\ell\gets$ their average
\EndFor
\State for each $\pi$, compute exact throttle/sharing factors $\phi_i(\pi),\psi_{ij}(\pi)$ and define $\widehat F(\pi):=F(\widehat{\mathbf d},\pi)$ via the structural model \eqref{eq:Cform} and recursion \eqref{eq:recurdef}
\State $\tilde\pi\gets\gamma$-oracle$\big(\min_{\pi\in\Pi}\widehat F(\pi)\big)$
\State \Return $\tilde\pi$
\end{algorithmic}
\end{algorithm}

\noindent\textbf{Correctness.} Lines 2--5 realize the structural-parameter
estimator (demands, rates, bandwidths) at the sample size of
Corollary~\ref{cor:uniform}, so the accuracy events \eqref{eq:Deps} and
\eqref{eq:rateacc} hold simultaneously at level $\varepsilon$ w.p.\ $\ge1-\delta$;
line~6 evaluates $\widehat F$ deterministically through the structural model
with the \emph{exact, configuration-specific} throttle factors (valid for
\emph{all} $\pi$ by Lemma~\ref{lem:propagate}, propagated accuracy
$\varepsilon_\star=\eta$), and line~7 returns a $\gamma$-approximate minimizer.
Corollary~\ref{cor:e2e} certifies the output guarantee. \textbf{Termination:}
finite loops, finite $\Pi$. \textbf{Complexity:} probing costs
$\sum_\ell r_\ell=O\!\big((m{+}1)\,Q'\,\frac{b_{\max}^2}{a_{\min}^2(\varepsilon-\beta)^2}\ln\frac{Q'}{\delta}\big)$
probes -- polynomial in the instance and \emph{independent of $|\Pi|$}; line~6
evaluates $\widehat F$ at the cost charged by the oracle's queries, and line~7
costs one call to the scheduling oracle. The reduction is black-box in the
optimizer, so the heterogeneous scheduling literature's approximation algorithms
plug in directly (subject to the model-matching caveat above for the achievable
$\gamma$).

\section{Summary of what is established}\label{sec:summary}

(a) Descriptors $L_i(\tau),N_{ij}(\tau)$ are defined and computed by
\textsc{ComputeDescriptors} (block-averaged measurements) and split into
schedule-independent structural constants and configuration-induced contention
state, the latter being an \emph{exactly computable} function of the schedule
(Definition~\ref{def:load}). (b) The estimator of the structural parameters
(demands, nominal rates, nominal bandwidths), propagated through the known
structural completion-time model \emph{evaluated with each configuration's own,
exactly-known contention throttles}, yields $\Chat$ and finish times that are
uniformly $(1\pm\eta)$-accurate \emph{over all configurations}
(Lemma~\ref{lem:propagate}), with explicit sample complexity (genuine Hoeffding
on independent blocks under finite-range dependence,
Proposition~\ref{prop:samples}) that honestly carries bias ($\beta$),
dependence range ($m$), and temporal drift ($\zeta_\ell H$,
Remark~\ref{rem:drift}). The load-dependence of completion time
(contention/queueing) is modeled, not assumed away: the realized
$C_u(\pi)$ depends on $\pi$ through $\phi_{x_u}(\pi)$, while degree-one
homogeneity in demand is preserved by the homogeneous-contention regime
(Assumption~\ref{ass:regime}); non-homogeneous nonlinearities are handled with
an explicit residual factor (Lemma~\ref{lem:lipschitz}). (c) The transfer
theorem (Theorems~\ref{thm:main}, \ref{thm:approxalg},
Corollary~\ref{cor:e2e}) shows that minimizing the estimated objective -- exactly
or $\gamma$-approximately -- yields a configuration within
$\gamma\frac{1+\eta}{1-\eta}$ of the true optimum for every
demand-scale-coherent objective (makespan under precedence, weighted completion
time, $\ell_p$ loads, $\dots$), and the constant $\frac{1+\eta}{1-\eta}$ is
tight. The $\gamma$ factor is inherited, with matching-model caveats, from the
scheduling literature; the estimation factor $\frac{1+\eta}{1-\eta}$ is the
contribution proved here, decoupled from optimization hardness.

\paragraph{Scope and honest limitations.} The headline guarantee is
\emph{conditional} on: (i) Assumption~\ref{ass:regime} (linear throughput with
multiplicative, demand-independent contention sharing) -- the standard
scheduling model, relaxed to a residual factor in Lemma~\ref{lem:lipschitz};
(ii) the contention state being a \emph{known} combinatorial function of the
schedule, so that throttles are computable without deploying $\pi$; (iii)
availability of a $\gamma$-approximation oracle for the (NP-hard) estimated
scheduling problem, whose constant depends on the machine/communication model
and is not universal for general DAG-with-communication scheduling; and (iv) the
probing model of Assumption~\ref{ass:probe} (bounded, $m$-dependent
work-counter probes with bounded relative bias). Within this scope the chain
measurement $\to$ estimation $\to$ mapping quality is fully rigorous and the
transfer constant is tight.
\end{document}
